% !TeX root = ../example.tex
\section{提示框}
\sectioncoverpage[assets/sigs_dom]{本节展示提示框样式与使用方法。} % 原图右边缘贴齐页面右边缘
\note{转入提示框：用有限的视觉强调区分建议、补充说明和警告。}
\subsection{样式演示}

\begin{frame}{提示、备注与警告}
      \begin{tipbox}[提示]
            使用 \texttt{tipbox} 放置操作建议，例如推荐的实验配置。
      \end{tipbox}

      \begin{notebox}[备注]
            使用 \texttt{notebox} 放置补充说明，例如数据来源或实验条件。
      \end{notebox}

      \begin{alertbox}[警告]
            使用 \texttt{alertbox} 强调易错点，例如复现时需要核对的设置。
      \end{alertbox}
\end{frame}
\note{比较提示、备注与警告的功能色。每个框只放一个需要强调的信息。}

% !TeX root = ../example.tex
\begin{frame}{内容与页面}
  \begin{columns}[T,onlytextwidth]
    \begin{column}{.46\textwidth}
      \textbf{大纲阶段}
      \begin{itemize}
        \item 组织章节与小节
        \item 补齐解释与证据
        \item 审阅并确认内容
      \end{itemize}
    \end{column}
    \begin{column}{.46\textwidth}
      \textbf{制作阶段}
      \begin{itemize}
        \item 划分页面与选择版式
        \item 分配正文与备注
        \item 检查并导出 PDF/PPTX
      \end{itemize}
    \end{column}
  \end{columns}
  \begin{notebox}[制作前提]
    详细大纲已确认；正文保留必要证据与限制。
  \end{notebox}
\end{frame}
\note{先对照左右两栏的大纲阶段与制作阶段，再指出下方“制作前提”说明框：它组织与主结构相关且必须可见的条件。条件用 notebox，操作建议用 tipbox，关键限制用 alertbox，定义和示例用对应环境。标题明确内容类别；不把每条正文都装箱，也不重复整页总结。}


\subsection{定义与示例}

\begin{frame}{定义与示例}
      \begin{definitionbox}[定义]
            \textbf{Modality}：信息的一种表示形式，如文本、图像、音频等。
            多模态学习关注如何联合建模这些异构来源。
      \end{definitionbox}

      \begin{examplebox}[示例]
            以 Tensor Fusion 为例，它通过外积显式保留单模态、双模态、
            三模态的交互，从而得到更丰富的联合表示。
      \end{examplebox}

\end{frame}
\note{比较定义与示例的语义。Tensor Fusion 文字沿用模板示例，正式论文讲解前应阅读并核实对应论文。}

\subsection{使用方法}

\begin{frame}[fragile]{提示框用法}
      \begin{verbatim}
\begin{tipbox}[自定义标题]
  需要强调的内容。
\end{tipbox}
      \end{verbatim}
      \begin{itemize}
            \item 将环境名替换为 \texttt{notebox}、\texttt{alertbox}、\texttt{examplebox} 或 \texttt{definitionbox}。
            \item 省略可选标题时，使用对应环境的默认标题。
      \end{itemize}
      \begin{tipbox}
            这是省略可选标题后的显示效果。
      \end{tipbox}
\end{frame}
\note{展示提示框环境的写法；省略标题会使用当前主题语言的默认名称。}
